% !TEX program = lualatex
\documentclass[11pt,a4paper]{article}
\usepackage[margin=2.2cm]{geometry}
\usepackage[italian]{babel}
\usepackage{alfdoc}
\usepackage{tikz-dimline}
\pgfplotsset{compat=newest}
\usepackage[colorlinks]{hyperref}
\title{\emo{seedling}\ Camminamenti e irrigazione dell'orto\\[0.3em]\large Pietre quadrate su ghiaia}
\author{Lezione per Alf, scritta da Codex e Claude}
\date{Mer 25 nov 2026}
% ---- materials: one colour each, the same in every figure ----
\definecolor{mSoil}{HTML}{8F7252}\definecolor{mBed}{HTML}{5E4632}\definecolor{mTile}{HTML}{B3B8B6}
\definecolor{mSand}{HTML}{DFC08A}\definecolor{mBase}{HTML}{A3A099}\definecolor{mChip}{HTML}{747A74}
\definecolor{mGravel}{HTML}{BDB3A2}\definecolor{mPebble}{HTML}{E2D8C8}\definecolor{mWood}{HTML}{A67B4C}
\definecolor{mString}{HTML}{D15C86}\definecolor{mLawn}{HTML}{A9BF7F}\definecolor{mRock}{HTML}{92989D}
\definecolor{mCloth}{HTML}{333638}\definecolor{mWater}{HTML}{3F8FC0}\definecolor{mLevel}{HTML}{D6C25A}
\tikzset{lab/.style={font=\small,align=center},mat/.style={fill=#1,draw=#1!70!black},
  wlab/.style={lab,text=white},
  pebble/.pic={\fill[mPebble] (0,0) ellipse [x radius=.075cm,y radius=.05cm];},
  chip/.pic={\fill[mChip] (-.08cm,-.06cm)--(.07cm,-.07cm)--(.09cm,.03cm)--(0,.09cm)--(-.08cm,.04cm)--cycle;},
  rock/.pic={\draw[mat=mRock] (-.24cm,-.1cm)--(-.15cm,.14cm)--(.06cm,.18cm)--(.24cm,.05cm)--(.17cm,-.14cm)--(-.05cm,-.18cm)--cycle;},
  leaf/.pic={\foreach \a in {0,72,...,288}{\fill[mLawn!80!black,rotate=\a] (0,0) to[bend left=45] (0,.26cm) to[bend left=45] (0,0);}},
  plan/.style={x=.52cm,y=.52cm},sect/.style={x=.25cm,y=.4cm}}
% \strew{x0}{y0}{x1}{y1}{nx}{ny}{pic}: a jittered grid of marks
\newcommand{\strew}[7]{\foreach \i in {1,...,#5}{\foreach \j in {1,...,#6}{
  \pgfmathsetmacro\px{(#1)+((#3)-(#1))*(\i-.5+.2*rand)/#5}\pgfmathsetmacro\py{(#2)+((#4)-(#2))*(\j-.5+.2*rand)/#6}
  \pgfmathsetmacro\pr{40*rand}\pic[rotate=\pr] at (\px,\py) {#7};}}}
% ---- plan: measured in cm, one drawing unit is 30.48 cm ----
\newcommand{\cm}[2]{\pgfmathsetmacro#1{#2/30.48}}
\cm\Ca{60}\cm\Cb{86}\cm\Cc{63}\cm\Cd{60} % four cross paths
\cm\La{122}\cm\Lb{366}\cm\Lc{122}        % three beds
\cm\W{122}\cm\Pw{60}                     % bed width, lawn path width
\pgfmathsetmacro\xA{\Ca}\pgfmathsetmacro\xB{\xA+\La}\pgfmathsetmacro\xC{\xB+\Cb}\pgfmathsetmacro\xD{\xC+\Lb}
\pgfmathsetmacro\xE{\xD+\Cc}\pgfmathsetmacro\xF{\xE+\Lc}\pgfmathsetmacro\xG{\xF+\Cd}\pgfmathsetmacro\yT{\Pw+\W}
\pgfmathsetmacro\cA{\Ca/2}\pgfmathsetmacro\cB{\xB+\Cb/2}\pgfmathsetmacro\cC{\xD+\Cc/2}\pgfmathsetmacro\cD{\xF+\Cd/2}
\newcommand{\beds}{\foreach \a/\b/\n in {\xA/\xB/1,\xC/\xD/2,\xE/\xF/3}{\draw[mat=mWood] (\a,\Pw) rectangle (\b,\yT);
  \fill[mBed] (\a+.25,\Pw+.25) rectangle (\b-.25,\yT-.25);\node[wlab,font=\footnotesize] at ({(\a+\b)/2},{\Pw+\W/2}) {Aiuola \n};}}
\newcommand{\paths}[1]{\fill[#1] (0,0) rectangle (\xG,\Pw);\foreach \a/\b in {0/\xA,\xB/\xC,\xD/\xE,\xF/\xG}{\fill[#1] (\a,\Pw) rectangle (\b,\yT);}}
\newcommand{\pathpebbles}{\strew{.2}{.1}{\xG-.2}{\Pw-.1}{70}{7}{pebble}
  \foreach \a/\b in {0/\xA,\xB/\xC,\xD/\xE,\xF/\xG}{\strew{\a+.15}{\Pw}{\b-.15}{\yT-.1}{7}{9}{pebble}}}
% lawn stones: one per cross path, one per small bed, three spaced evenly between the middle cross paths
\pgfmathsetmacro\mA{(\xA+\xB)/2}\pgfmathsetmacro\mB{(\cB+\cC)/2}\pgfmathsetmacro\mC{(\xE+\xF)/2}
\newcommand{\tiles}{\foreach \tx in {\cA,\mA,\cB,{\cB+(\cC-\cB)/4},\mB,{\cB+3*(\cC-\cB)/4},\cC,\mC,\cD}{
    \draw[mat=mTile,rounded corners=.5pt] ({\tx-.667},\Pw/2-.667) rectangle +(1.333,1.333);}
  \foreach \c in {\cA,\cB,\cC,\cD}{\draw[mat=mTile,rounded corners=.5pt] (\c-.667,\Pw+\W/2-.667) rectangle +(1.333,1.333);}}
% ---- section, in units of 2.5 cm: tile top at 0, flat bed at -3, 2.5 cm of gravel, 1 cm of sand, 4 cm of tile ----
\newcommand{\soil}{\fill[mSoil] (0,-3) rectangle (60,-6.5);\node[wlab] at (30,-5) {Terra compattata};}
\newcommand{\cloth}{\draw[mCloth,line width=1.5pt] (0,-3)--(60,-3);}
\newcommand{\base}{\fill[mBase] (0,-3) rectangle (60,-2);\strew{0.3}{-2.95}{59.7}{-2.05}{54}{2}{chip}}
\newcommand{\sand}{\foreach \x in {10,34}{\fill[mSand] (\x,-2) rectangle +(16,.4);}}
\newcommand{\sectiles}{\foreach \x in {10,34}{\draw[mat=mTile,rounded corners=1pt] (\x,-1.6) rectangle +(16,1.6);}}
\newcommand{\level}[2]{\draw[mat=mLevel,rounded corners=1pt] ([xshift=-.8cm]#1,#2) rectangle +(1.6cm,.3cm);\draw[fill=mLawn!50,draw=mLevel!60!black,rounded corners=1pt] ([xshift=-.3cm,yshift=.08cm]#1,#2) rectangle +(.6cm,.14cm);}
\setlength{\parindent}{0pt}\setlength{\parskip}{0.5em}
\begin{document}
\maketitle
\pgfmathsetseed{7}
\section*{Il risultato}
Il camminamento passa lungo il prato e attraversa l'orto in quattro traversi. Usiamo tredici piastrelle quadrate da 40 cm (16 in). Sul lato prato ne mettiamo una davanti a ogni traverso, una al centro di ogni aiuola piccola, e tre lungo l'aiuola grande. Le altre quattro stanno al centro dei traversi. Sotto c'è un fondo piatto con 2,5 cm di ghiaia, e sotto ogni piastrella uno strato di sabbia. Le piastrelle restano più alte della ghiaia, quindi camminiamo sulle pietre. Il lato della recinzione, con le rocce, resta per una fase futura. Il lato prato e i due traversi esterni sono larghi 60 cm. I due traversi interni sono larghi 86 e 63 cm. La fila intera è lunga 8,8 m.
\par\nopagebreak
\begin{center}\begin{tikzpicture}[plan]
\fill[mLawn] (-.6,-.9) rectangle (\xG+.6,0);\fill[mSoil] (-.6,0) rectangle (\xG+.6,\yT+1.1);
\paths{mGravel}\pathpebbles\beds\tiles
\foreach \a/\b in {\xA/\xB,\xC/\xD,\xE/\xF}{\foreach \t in {.2,.5,.8}{\pic at ({\a+\t*(\b-\a)},\Pw+.95) {leaf};\pic at ({\a+\t*(\b-\a)},\yT-.95) {leaf};}}
\foreach \i in {0,...,20}{\pgfmathsetmacro\rx{\xA+\i*(\xF-\xA)/20}\pgfmathsetmacro\rr{180*rnd}\pic[rotate=\rr] at (\rx,\yT+.5) {rock};}
\draw[mat=mWood] (-.6,\yT+1.1) rectangle (\xG+.6,\yT+1.4);
\node[lab] at (\xG/2,-.45) {Prato};\node[lab,above] at (\xG/2,\yT+1.4) {Recinzione};
\node[wlab] at (\xG-.6,\yT+.6) {Rocce};
% dimensions, light and tiny, with tikz-dimline
\colorlet{mDim}{black!25}
\newcommand{\dm}[5][90]{\dimline[color=mDim,extension style={draw=mDim,line width=.4pt},extension start length=0cm,extension end length=0cm,extension start angle=#1,extension end angle={-(#1)},
  label style={fill=none,font=\scriptsize,inner sep=1.5pt,text=black!40,#5}]{#2}{#3}{#4}}
% construction lines, from the beds' corners and the path's ends to the dimensions
\begin{scope}[every path/.style={draw opacity=.3}]
\foreach \x in {\xA,\xB,\xC,\xD,\xE,\xF}{\draw (\x,\Pw)--(\x,-1.65);}
\foreach \x in {0,\xG}{\draw (\x,0)--(\x,-2.5);}
\foreach \y in {\Pw,\yT}{\draw (\xA,\y)--(-1,\y);}
\draw (0,0)--(-1,0);
\end{scope}
\foreach \a/\b/\t in {0/\xA/60,\xA/\xB/120,\xB/\xC/86,\xC/\xD/370,\xD/\xE/63,\xE/\xF/120,\xF/\xG/60}{\dm{(\a,-1.65)}{(\b,-1.65)}{\t}{above}}
\dm{(0,-2.5)}{(\xG,-2.5)}{8,8 m}{above}
\node[above left,font=\scriptsize,inner sep=1.5pt,text=black!40] at (\xG,-2.5) {Misure in cm};
\dm[-90]{(-1,0)}{(-1,\Pw)}{60}{sloped=false,left}\dm[-90]{(-1,\Pw)}{(-1,\yT)}{120}{sloped=false,left}
\end{tikzpicture}\end{center}

\section{Tracciare e scavare \emo{check mark button}}
Abbiamo teso lo spago sui due bordi e scavato il vialetto largo 60 cm. Lo spago segna solo i bordi, quindi può seguire la pendenza del terreno. Prima di posare proviamo il passo con le piastrelle, perché tra una pietra e l'altra restano 50--70 cm.
\par\nopagebreak
\begin{center}\begin{tikzpicture}[plan]
\paths{mSoil!35}\beds
\draw[mString,line width=1.2pt] (0,0) rectangle (\xG,\Pw);
\foreach \a/\b in {0/\xA,\xB/\xC,\xD/\xE,\xF/\xG}{\draw[mString,line width=1.2pt] (\a,\Pw) rectangle (\b,\yT);
  \foreach \x in {\a,\b}{\foreach \y in {\Pw,\yT}{\draw[mat=mWood] (\x-.2,\y-.2) rectangle +(.4,.4);}}}
\foreach \i in {0,...,4}{\draw[mat=mWood] (\i*\xG/4-.2,-.2) rectangle +(.4,.4);}
\tiles
\draw[mat=mWood] (8,-1.6) rectangle +(.4,.4);\node[lab,right] at (8.5,-1.4) {Picchetti};
\draw[mString,line width=1.2pt] (16,-1.4)--(18,-1.4);\node[lab,right] at (18.1,-1.4) {Spago};
\end{tikzpicture}\end{center}

\section{Spianare il fondo}
Togliamo l'erba e la terra soffice rimaste. Nel senso della lunghezza seguiamo il terreno, togliendo buche e gobbe. Di traverso lasciamo una leggera pendenza dal prato verso la scarpata. Poi versiamo un secchio d'acqua e guardiamo dove va. Se si ferma in una pozza, o contro il legno delle aiuole, correggiamo il fondo in quel punto. Le rocce restano dove sono, perché possono sostenere il terreno delle aiuole.
\par\nopagebreak
\begin{center}\begin{tikzpicture}[x=.2cm,y=.2cm]
% across the path, in units of 2.5 cm; heights illustrative
\fill[mSoil] (-12,10) to[out=-5,in=125] (0,0)--(24,-1.2)--(40,-2)--(40,-6)--(-12,-6)--cycle;
\draw[line width=2pt,mLawn] (-12,10) to[out=-5,in=125] (0,0);
\draw[mat=mWood] (24,-2) rectangle (40,6.8);\fill[mBed] (25.5,-1.8) rectangle (40,5.8);
\node[wlab] at (32.5,2) {Aiuola};
\draw[mat=mTile] (7,6)--(11,6)--(10.4,2.5)--(7.6,2.5)--cycle;
\node[lab,above] at (9,6) {Secchio d'acqua};
\draw[mWater,line width=1.5pt,->] (2,.9)--(20,0);
\node[lab,above] at (-9,10.5) {Prato};
\draw[<->] (0,-3.2)--(24,-3.2) node[wlab,midway,below]{60 cm};
\end{tikzpicture}\end{center}

\section{Compattare e stendere il telo}
Compattiamo il terreno leggermente umido con il battitore (\textit{hand tamper}). Il fondo porta tutto il peso, quindi lo battiamo bene. Poi stendiamo il geotessile (\textit{landscape fabric}), un telo permeabile che tiene separate la terra e la ghiaia. Dove finisce un pezzo di telo e ne comincia un altro, il secondo copre il primo per almeno 15 cm, così la terra non passa dalla fessura. Fermiamo i bordi con picchetti metallici a U.
\par\nopagebreak
\begin{center}\begin{tikzpicture}[sect]
\soil\cloth
\draw[mat=mWood] (41.6,-2) rectangle (42.4,3);\draw[mat=mRock] (38.5,-2.8) rectangle (45.5,-2);
\node[lab,right] at (42.6,2.2) {Battitore};\draw[->] (18,-.6) node[lab,above]{Geotessile}--(18,-2.8);
\end{tikzpicture}\end{center}

\section{Posare i tubi}
Un tubo unico arriva dal lato destro dell'orto e corre sul telo nella striscia di 10 cm tra il legno delle aiuole e le piastrelle, dove la ghiaia lo copre. Davanti a ogni aiuola una biforcazione a T sale, scavalca il bordo di legno e, con la sua valvola, alimenta i gocciolatori. Il tubo resta fuori dalle piastrelle, che altrimenti dondolerebbero, e non ci battiamo sopra con il battitore. L'antiriflusso impedisce all'acqua dell'impianto di tornare nella rete potabile, e il filtro e il regolatore di pressione proteggono i gocciolatori \cite{ua}. Apriamo l'acqua e controlliamo i raccordi prima di coprire il tubo.
\par\nopagebreak
\begin{center}\begin{tikzpicture}[plan,x=.5cm]
% from above, lawn at the bottom
\paths{mGravel!60}\beds
\foreach \tx in {\cA,\mA,\cB,{\cB+(\cC-\cB)/4},\mB,{\cB+3*(\cC-\cB)/4},\cC,\mC,\cD}{\draw[mat=mTile,rounded corners=.5pt] ({\tx-.667},\Pw/2-.667) rectangle +(1.333,1.333);}
\foreach \c in {\cA,\cB,\cC,\cD}{\draw[mat=mTile,rounded corners=.5pt] (\c-.667,\Pw+\W/2-.667) rectangle +(1.333,1.333);}
\pgfmathsetmacro\yM{\Pw-.16}
\draw[mWater,line width=1.5pt] (\xG+2.4,\yM)--(\mA,\yM);
\foreach \a/\b in {\xA/\xB,\xC/\xD,\xE/\xF}{\pgfmathsetmacro\m{(\a+\b)/2}
  \draw[mWater,line width=1.5pt] (\m,\yM)--(\m,\Pw+.7);
  \draw[mWater,line width=1pt] (\a+.6,\Pw+.7)--(\b-.6,\Pw+.7);
  \foreach \x in {\a+.6,\b-.6}{\draw[mWater,line width=1pt] (\x,\Pw+.7)--(\x,\yT-.6);\foreach \q in {.9,1.7,2.5}{\fill[white] (\x,\yT-\q) circle [radius=.05cm];}}
  \draw[mat=mString] (\m-.2,\Pw+.3) rectangle +(.4,.4);}
\draw[mat=mLevel] (\xG+.4,\yM-.35) rectangle (\xG+1.9,\yM+.35);\draw[mat=mWater] (\xG+2.2,\yM-.35) rectangle (\xG+2.9,\yM+.35);
\node[lab,anchor=north east] at (\xG+2.9,-.4) {Rubinetto, antiriflusso,\\filtro, regolatore, timer};
\draw[->] (\xG+1.6,-.4)--(\xG+1.6,\yM-.4);
\draw[->] (9.6,-.7) node[lab,below]{Tubo nella ghiaia}--(9.6,\yM-.1);
\draw[->] (16.4,-.7) node[lab,below]{Valvola}--(\mB+.95,\Pw+.1)--(\mB+.25,\Pw+.45);
\end{tikzpicture}\end{center}

Ecco lo stesso punto visto dal lato destro, con le spalle al rubinetto e l'aiuola 3 davanti a noi.
\par\nopagebreak
\begin{center}\begin{tikzpicture}[x=.1cm,y=.1cm]
% in cm, to scale except the height of the lawn; the bed's 20 cm is an estimate
\fill[mSoil] (-28,28) to[out=-5,in=125] (0,2.5)--(0,0)--(118,0)--(118,-10)--(-28,-10)--cycle;
\draw[line width=2pt,mLawn] (-28,28) to[out=-5,in=125] (0,2.5);
\fill[mBase] (0,0) rectangle (60,2.5);\strew{0.5}{.3}{59.5}{2.2}{34}{1}{chip}
\draw[mCloth,line width=1.5pt] (0,0)--(60,0);
\fill[mSand] (10,2.5) rectangle (50,3.5);
\draw[mat=mTile,rounded corners=1pt] (10,3.5) rectangle (50,7.5);\node[lab] at (30,5.5) {Piastrella};
\draw[mat=mWood] (60,-2) rectangle (64,20);\fill[mBed] (64,-2) rectangle (118,17.5);
\node[wlab] at (95,7) {Aiuola 3};
% the pipe in the gravel, its branch over the wood, and the drip line
\draw[mWater,line width=2pt] (56,1.3)--(56,23)--(68,23)--(68,18.7)--(112,18.7);
\fill[mWater,draw=white] (56,1.3) circle [radius=.1cm];
\draw[mat=mString] (66,20) rectangle (70,22.5);
\foreach \x in {78,90,102}{\fill[white,draw=mWater] (\x,18.7) circle [radius=.07cm];}
\node[lab,above] at (-18,29) {Prato};
\draw[->] (40,17) node[lab,above]{Tubo nella ghiaia}--(54.6,2.3);
\draw[->] (84,27) node[lab,right]{Valvola}--(70.5,21.8);
\draw[->] (5,14) node[lab,above]{Ghiaia}--(5,2.8);
\foreach \a/\b/\t in {0/10/10,10/50/40,50/60/10}{\draw[white,<->] (\a,-4)--(\b,-4) node[wlab,midway,below,font=\scriptsize]{\t};}
\node[wlab,right,font=\scriptsize] at (61,-5.8) {Misure in cm};
\end{tikzpicture}\end{center}

\section{Stendere la ghiaia}
La ghiaia è il \textit{paver base} che abbiamo già, roccia frantumata che si incastra quando la compattiamo. Ne stendiamo 2,5 cm su tutto il camminamento, la livelliamo con il rastrello, e la battiamo. È uno strato sottile, adatto a camminare su pietre distanziate e terreno fermo. Le pavimentazioni continue richiedono una base più spessa \cite{cmha}.
\par\nopagebreak
\begin{center}\begin{tikzpicture}[sect]
\soil\base\cloth
\draw[<->] (61,-3)--(61,-2) node[lab,midway,right]{2,5 cm};
\draw[->] (30,-.6) node[lab,above]{Ghiaia}--(30,-1.9);
\end{tikzpicture}\end{center}

\section{Posare le piastrelle}
Sotto ogni piastrella stendiamo 1 cm di sabbia (\textit{paver sand}), che la sostiene in modo uniforme. Cominciamo con due piastrelle. Se una dondola quando premiamo un angolo, la solleviamo e correggiamo la sabbia. La livella appoggiata sopra deve mostrare il lato verso la scarpata più basso di 4--8 mm. Quando vanno bene, posiamo le altre undici allo stesso modo.
\par\nopagebreak
\begin{center}\begin{tikzpicture}[sect]
\soil\base\cloth\sand\sectiles\level{42}{.05}\node[lab,above] at (42,.9) {Livella};
\node[lab] at (18,-.75) {Piastrella};
\draw[->] (4,.6) node[lab,above]{Sabbia}--(10.6,-1.75);
\draw[->] (30,.6) node[lab,above]{Ghiaia}--(30,-1.9);
\draw[<->] (61,0)--(61,-3) node[lab,midway,right]{7,5 cm};
\end{tikzpicture}\end{center}

\section{Controllare dopo la pioggia}
Il prato è molto più in alto dell'orto e scende con una pendenza di 45° o più, mentre la scarpata è più in basso. L'acqua scende dal prato, attraversa il camminamento, e prosegue nei traversi verso la scarpata. Controlliamo che il legno delle aiuole non fermi l'acqua, che la ghiaia resti pulita dalla terra del prato, e che la sabbia rimanga sotto le piastrelle. Poi premiamo gli angoli di ogni pietra. Dove una pietra si muove, la solleviamo, correggiamo la sabbia, e la rimettiamo.
\par\nopagebreak
\begin{center}\begin{tikzpicture}[x=.125cm,y=.125cm]
% across the garden, in units of 2.5 cm, lawn on the left; widths of path and bed to scale, heights illustrative
\fill[mSoil] (0,17)..controls (11,15.5) and (17,9)..(24,0)--(24,-3)--(48,-3)--(48,-1)--(97,-2)--(108,-16)--(120,-16)--(120,-22)--(0,-22)--cycle;
\fill[mLawn] (0,17)..controls (11,15.5) and (17,9)..(24,0)--(24,1.3)..controls (17.8,10.2) and (11.5,16.9)..(0,18.3)--cycle;
\begin{scope}[shift={(24,0)},yslant=-.02,shift={(-24,0)}]
\fill[mBase] (24,-3) rectangle (48,-2);\fill[mSand] (28,-2) rectangle (44,-1.6);
\draw[mCloth,line width=1.5pt] (24,-3)--(48,-3);
\draw[mat=mTile,rounded corners=.5pt] (28,-1.6) rectangle (44,0);
\draw[mat=mWood] (48,-2) rectangle (96,7);\fill[mBed] (50,-2) rectangle (94,6);
\node[wlab] at (72,2) {Aiuola};
\end{scope}
\foreach \x/\y/\a in {98/-2/0,101.5/-6/60,105/-10.5/20,108/-14.5/90}{\pic[rotate=\a] at (\x,\y) {rock};}
\draw[mat=mWood] (116,-20) rectangle (118,10);\node[lab,above] at (116,10) {Recinzione};
\node[lab,above] at (7,17.6) {Prato};\node[lab,above] at (36,1) {Camminamento};
\node[lab,right] at (101,-1) {Rocce};
\node[wlab] at (60,-14) {Terra compattata};\node[wlab,left] at (104,-19) {Scarpata};
\draw[mWater,line width=1.2pt,->] (17,13)--(25,2.7);
\draw[mWater,line width=1.2pt,->] (103,-2.8)--(110.5,-12.5);
\node[lab,right] at (21,10) {Deflusso};
\foreach \x in {21,30,...,111}{\draw[mWater,line width=1pt] (\x,20)--+(-1,-2.4) (\x+4.5,17.5)--+(-1,-2.4);}
\end{tikzpicture}\end{center}
\section*{Quanto materiale serve}
Il camminamento copre 8,6 m². I sacchi sono da 14 L (0,5 ft³), arrotondati per eccesso. Per la ghiaia contiamo una resa dell'85\% dopo la compattazione, che è un'ipotesi. La ghiaia che abbiamo basta appena. Della sabbia teniamo 3 sacchi e ne cambiamo 5.
\begin{center}\small
\begin{tabularx}{\linewidth}{@{}Xrrr@{}}
\toprule Materiale & Volume posato (L) & Sacchi necessari & Sacchi comprati\\ \midrule
Ghiaia, 2,5 cm su tutto il camminamento & 210 & 18 & 18\\
Sabbia, 1 cm sotto le piastrelle & 21 & 2 & 8\\
\bottomrule
\end{tabularx}\end{center}
\begin{thebibliography}{CMHA}
\bibitem[CMHA]{cmha} Concrete Masonry \& Hardscapes Association, \textit{Construction of interlocking concrete pavements}, PAV-TEC-002. \url{https://www.cmha.org/resource/pav-tec-002/}.
\bibitem[UA]{ua} University of Arizona Cooperative Extension, \textit{Drip irrigation tune-up}. \url{https://extension.arizona.edu/sites/default/files/2024-10/DripIrrigationTune-up_0.pdf}.
\end{thebibliography}
\revisiontable
\end{document}
